\begin{abstract}
从一段由多个镜头组成的单目视频中，在统一世界坐标系下重建人体运动是一项具有挑战性的任务，因为突然发生的镜头切换会带来大幅视角变化，并显著改变画面中的可见场景和人体外观。现有在线方法通常假设输入在时间上连续；已有多镜头方法则联合处理镜头或优化已观测序列，而不是在帧到达时递归地产生预测。在镜头切换处，递归历史为恢复镜头间世界坐标关系提供参照，但将同一状态继续传播到新镜头可能引起随时间变化的相机漂移，而单独重新初始化又会丢失这一参照。为处理这一冲突，我们提出 \method{}，一个通过解耦镜头间对齐与递归状态传播，实现多镜头视频流式多人体四维重建的在线框架。在每个镜头边界，历史条件路径估计镜头间变换，独立重新初始化的路径则提供新镜头中继续传播的状态。随后，基于几何的关联建立跨镜头身份对应，共享变换将相机和人体配准到已有世界坐标系。在三个多人体基准上，相比跨镜头直接延续递归状态，\method{} 将按数据集等权计算的平均 W-MPJPE 降低 9.5\%，并将平均 IDF1 绝对提高 0.099。在 EgoBody 预定义的极端视角子集上，\method{} 将 W-MPJPE 降低 29.9\%，同时保持无需未来帧或逐视频优化的逐帧流式推理。
\end{abstract}
